% !TeX root = ../../金融经济学_王江_讲义及习题.tex
% 知识点：references/金融经济学/金融经济学-王江_261013.pdf，印刷页 49–62
% 补充：references/金融经济学/金融经济学_第4章 套利和资产定价.pptx（第 1–36 页）
% 习题：references/金融经济学/史上最全最完整的金融经济学答案(王江版).pdf，页 14–18（印刷页 13–17）

\lecture{套利和资产定价}

\ifdefined\WJHandoutKnowledge

第 3 章在 Arrow--Debreu 市场结构下讨论均衡与效率；本章回到第 2 章的一般支付矩阵 $X$，讨论\textbf{证券价格的基本性质}、\textbf{定价原理}、\textbf{无套利}与\textbf{资产定价基本定理}。

\begin{tcolorbox}[knowledge,title=4.1 一般市场结构]
多数可交易证券在多个状态下都有正支付，称为\textbf{复合证券}（相对只在单一状态付 1 的阿罗证券而言）。设 1 期有 $\Omega$ 个状态，市场上有 $N$ 只证券（$n=1,\dots,N$），第 $n$ 只证券的支付向量为 $X_n=\begin{bmatrix}x_{1,n} \\ \vdots \\ x_{\Omega,n}\end{bmatrix}$，\textbf{市场结构矩阵}为（与第 2.3 节一致，行为状态、列为证券）：
\[
    X=[X_{\cdot,1},\,\cdots,\,X_{\cdot,N}]=
    \begin{bmatrix}
        x_{1,1}   & \cdots & x_{1,N}   \\
        \vdots    & \ddots & \vdots    \\
        x_{\omega,1} & \cdots & x_{\omega,N} \\
        \vdots    & \ddots & \vdots    \\
        x_{\Omega,1} & \cdots & x_{\Omega,N}
    \end{bmatrix}.
\]

\tcbline

\textbf{A.\ 冗余证券}\quad 若证券 $j$ 的支付列 $x_j$ 可由其余证券的\textbf{复制组合}（replicating portfolio）线性表出，即存在 $\theta_{\setminus j}$ 使得
\[
    x_j=X_{\setminus j}\theta_{\setminus j},
\]
则称 $j$ 为\textbf{支付意义上的冗余证券}。对任意原组合 $\theta$，
\[
 X\theta=X_{\setminus j}\theta_{\setminus j}+\theta_jx_j
 =X_{\setminus j}(\theta_{\setminus j}+\theta_j\theta^*_{\setminus j}),
\]
其中 $\theta^*_{\setminus j}$ 是复制 $x_j$ 的固定配方。这说明删去冗余列不改变可达支付集合。若市场无摩擦且满足一价定律，还必须有
\[
 S_j=S_{\setminus j}^{\mathsf T}\theta^*_{\setminus j},
\]
所以替换前后的成本也相同；若价格不相等，同支付的两种实现方式反而会产生套利。资源配置上的等价还要求复制交易不受限制，原有证券禀赋也按复制关系换算。存在交易成本、卖空约束等摩擦时，支付可复制不保证交易上可替代。

\begin{tcolorbox}[example,title=PPT 补充：第三只证券是否增加支付维度]
晴、雨两个状态下，证券 $A,B,C$ 的支付为
\[
 X=\begin{bmatrix}10&0&20\\0&10&10\end{bmatrix}.
\]
$C=2A+B$，故第三列虽不同于前两列，却没有增加独立支付维度；无套利价格必须满足 $S_C=2S_A+S_B$。例如原来持有 $aA+bB+cC$，可改为 $(a+2c)A+(b+c)B$，各状态支付与成本均相同。\textbf{证券种类数增加，不等于市场更完全。}
\end{tcolorbox}

\tcbline

\textbf{B.\ 证券市场的等价描述}\quad 无摩擦时，可用\textbf{线性无关}的 $N$ 只证券描述市场，$\mathrm{rank}(X)=N$（通常设 $N\le\Omega$）。任取 $N$ 个线性无关的组合向量 $\theta_1,\dots,\theta_N$，令 $H=[\theta_1,\cdots,\theta_N]$（$N\times N$ 可逆），则新支付矩阵 $X_\theta=XH$ 与原市场\textbf{等价}：原证券列空间不变，$X=X_\theta H^{-1}$，每只原证券都可由新组合复制。更一般地，原组合 $\theta$ 对应新组合 $H^{-1}\theta$，因为
\[
 X\theta=XH(H^{-1}\theta)=X_\theta(H^{-1}\theta).
\]
新证券的价格须同时变为 $S_\theta=H^{\mathsf T}S$，于是
\[
 S_\theta^{\mathsf T}H^{-1}\theta=S^{\mathsf T}\theta.
\]
支付和成本都不变，这才是同一市场的等价描述。可逆变换只是更换支付空间的一组基；若 $H$ 不可逆，可能丢掉支付维度。秩不变也可由 $X=XHH^{-1}$ 得到
\[
 \mathrm{rank}(X)\le\mathrm{rank}(XH)\le\mathrm{rank}(X),
 \qquad\mathrm{rank}(XH)=N.
\]

\tcbline

\textbf{C.\ 生成（span）与完全市场}\quad 记 $\mathbf{1}_\omega$ 为第 $\omega$ 分量为 1、其余为 0 的列向量（状态 $\omega$ 阿罗证券的支付）。若 $\mathrm{rank}(X)=\Omega$ 且 $N\ge\Omega$，则对每一 $\omega$ 存在组合 $\theta_\omega$ 使 $X\theta_\omega=\mathbf{1}_\omega$。当 $X$ 为 $\Omega\times\Omega$ 可逆方阵时，解唯一：
\[
    \theta_\omega=X^{-1}\mathbf{1}_\omega.
\]
此时称不确定性由证券\textbf{生成}（span）；市场与第 3 章完全阿罗市场等价。

\begin{tcolorbox}[knowledge]
\textbf{定理 4.1}\quad 无摩擦时，证券市场\textbf{完全}当且仅当线性无关证券只数等于状态数（可取 $N=\Omega$ 且 $\mathrm{rank}(X)=\Omega$，此时 $X$ 可逆）。这与定义 3.2、\textcolor{red}{\textbf{行满秩}}判据一致。
\end{tcolorbox}

\begin{tcolorbox}[example,title=PPT 补充：先复制阿罗证券，再复制任意支付]
两状态、两证券的支付矩阵为
\[
 X=\begin{bmatrix}2&3\\5&1\end{bmatrix},\qquad
 X^{-1}=\frac1{13}\begin{bmatrix}-1&3\\5&-2\end{bmatrix}.
\]
因此 $\theta_1=[-1,5]^{\mathsf T}/13$、$\theta_2=[3,-2]^{\mathsf T}/13$ 分别复制 $[1,0]^{\mathsf T}$、$[0,1]^{\mathsf T}$。要复制支付 $y=[5,2]^{\mathsf T}$，先把它写成五份第一状态阿罗证券加两份第二状态阿罗证券，再换成可交易证券：
\[
 \theta=X^{-1}y=5\theta_1+2\theta_2
 =\frac1{13}\begin{bmatrix}1\\21\end{bmatrix},\qquad
 X\theta=\frac1{13}\begin{bmatrix}2+63\\5+21\end{bmatrix}
 =\begin{bmatrix}5\\2\end{bmatrix}.
\]
复制成本为 $(S_1+21S_2)/13$。阿罗证券不必实际挂牌交易，能用已有证券复制即可。
\end{tcolorbox}
\end{tcolorbox}

\begin{tcolorbox}[knowledge,title=4.2 套利]
记交易证券的 0 期价格向量为 $S=\begin{bmatrix}S_1 \\ \vdots \\ S_N\end{bmatrix}$。组合 $\theta=\begin{bmatrix}\theta_1 \\ \vdots \\ \theta_N\end{bmatrix}$ 的 0 期成本为 $S^{\mathsf T}\theta$，1 期支付为 $X\theta$。从支付 $X$ 到价格 $S$ 的映射称为\textbf{资产定价关系}（资产定价模型）。若 $X\theta\ge 0$，称该组合在 1 期满足\textbf{有限责任}。

\begin{tcolorbox}[knowledge]
\textbf{定义 4.1}\quad 组合 $\theta$ 构成\textbf{套利机会}，若同时满足：
\begin{enumerate}
    \item $S^{\mathsf T}\theta\le 0$（0 期净支出非正）；
    \item $X\theta\ge 0$（1 期各状态支付非负）；
    \item 上述两式中\textbf{至少有一个严格不等式}成立（0 期严格收入或某状态严格正支付）。
\end{enumerate}
等价地，扩展支付向量 $\begin{bmatrix}-S^{\mathsf T}\theta \\ X\theta\end{bmatrix}$ 非负且至少一分量严格为正。
\end{tcolorbox}

\textbf{向量不等号约定：}$z\ge0$ 指各分量非负，$z>0$ 指 $z\ge0$ 且 $z\ne0$（至少一个分量为正），$z\gg0$ 才指每个分量都严格为正。因此套利不要求每个未来状态都获利，只要求没有亏损并在某时点、某状态有正所得。

常分为三类（记严格不等方向与教材一致）：
\begin{itemize}
    \item \textbf{第 1 类套利：}$S^{\mathsf T}\theta<0$ 且 $X\theta=0$——0 期白捡钱、未来无负债；
    \item \textbf{第 2 类套利：}$S^{\mathsf T}\theta=0$ 且 $X\theta>0$——\textbf{套利组合}（零成本、未来非负且非零支付）；
    \item \textbf{第 3 类套利：}$S^{\mathsf T}\theta<0$ 且 $X\theta>0$——0 期收入为正且未来非负、非零。
\end{itemize}

初始成本为零时，买入恰由卖空收入融资；初始成本为负时，卖空收入支付买入成本后还有余额。这就是本章“不需要额外初始资源”的含义。PPT 把前者称作自融资；动态模型中的自融资还要求各期调仓不追加外部资金，不能仅凭初始成本为零就判断。

\begin{tcolorbox}[example]
\textbf{例 4.1}\quad 三只证券的价格写在树的根上，支付写在各枝：
\begin{center}
\begin{tabular}{@{}c@{\quad}c@{\quad}c@{}}
\textbf{A} & \textbf{B} & \textbf{C} \\[0.35em]
\WJEventTree[left={$1$},stem]{1,1,1} &
\WJEventTree[left={$1$},stem]{0,2,2} &
\WJEventTree[left={$2$},stem]{2,0,0}
\end{tabular}
\end{center}
即 $S^{\mathsf T}=[1,\,1,\,2]$，$X=\begin{bmatrix}1&0&2\\1&2&0\\1&2&0\end{bmatrix}$。

组合 $\theta_1=[2;\,-1;\,-1]$：买入 $2$ 单位 $A$、卖出 $1$ 单位 $B$ 和 $1$ 单位 $C$。$S^{\mathsf T}\theta_1=-1$，$X\theta_1=0$，为\textbf{第 1 类套利}。支付 $[2;2;2]$ 可由 $2$ 单位 $A$ 以成本 $2$ 复制，也可由 $1$ 单位 $B$ 加 $1$ 单位 $C$ 以成本 $3$ 复制——同支付不同价。

另一种看法是先比较 $B+C$ 与 $2A$：二者都支付 $[2,2,2]^{\mathsf T}$，但成本分别为 $3$ 和 $2$，故应卖出昂贵组合、买入便宜组合。

组合 $\theta_2=[2;\,0;\,-1]$：买入 $2$ 单位 $A$、卖出 $1$ 单位 $C$。$S^{\mathsf T}\theta_2=0$，$X\theta_2=[0;\,2;\,2]>0$，为\textbf{第 2 类套利}。
\end{tcolorbox}
\end{tcolorbox}

\begin{tcolorbox}[knowledge,title=4.3 无套利原理]
套利机会只依赖公开的支付 $X$ 与价格 $S$，不依赖私人信息或对状态概率的主观信念。

\begin{tcolorbox}[knowledge]
\textbf{定理 4.2}\quad 市场均衡中\textbf{不存在}套利机会。
\end{tcolorbox}

\textbf{证明（按 PPT 的消费增量写法）：}反设均衡价格下存在套利 $\theta$。给任一参与者原有证券头寸再加上 $\theta$，同时把初始净收入用于 0 期消费，消费可由 $c_k$ 提高为
\[
 c'_k=c_k+\begin{bmatrix}-S^{\mathsf T}\theta\\X\theta\end{bmatrix}
 \ge c_k,\qquad c'_k\ne c_k.
\]
该交易的初始成本非正，不需要追加初始财富；未来也没有负支付。若至少一名参与者对这种非负、非零的消费增加严格偏好，则 $U_k(c'_k)>U_k(c_k)$，与均衡时其消费最优矛盾。\textbf{推出矛盾只需一次可行改进，无须先诉诸无限放大。}无限放大是在无交易限制下说明原价格无法维持的直觉。

\textbf{前提与边界：}市场无摩擦，可按给定价格买卖并实施所需空头；偏好须对上述消费增量具有严格单调性（本章不满足性假设用于此处）。仅有“某处还想多消费”的局部不满足性，并不足以保证对任意套利支付都严格偏好。无套利也不同于信息有效：这里检验的是价格与逐状态支付是否相容，不要求价格已经反映全部信息。均衡 $\Rightarrow$ 无套利是必要条件；无套利 $\not\Rightarrow$ 均衡，还需要偏好最优化与市场出清。\textbf{无摩擦是交易制度假设，无套利是给定支付与价格的约束，两者不等价}：例 4.1 即可在无摩擦条件下呈现错误价格与套利。

\begin{tcolorbox}[example,title=PPT 补充：跨市场的一价定律]
同一时刻、同一可交割黄金在两个市场分别报价 $1900$ 与 $1905$ 美元。假设可同时成交、自由卖空和交割，且没有交易与运输成本，买入便宜的黄金、卖出昂贵的黄金可立即收入 $5$ 美元，净实物头寸为零，是第 1 类套利。买入压力与卖出压力解释价差如何趋于消失，但不能据此推出最终价格必为两报价的中点。有摩擦时须扣除可锁定的全部成本；承担尚未对冲的汇率或交割风险的交易不能仅凭预期价差称为套利。这里数字只是教学假设。
\end{tcolorbox}

\begin{tcolorbox}[knowledge]
\textbf{定义 4.2}\quad 证券市场上不存在套利机会，称为满足\textbf{无套利原理}（no-arbitrage principle）。
\end{tcolorbox}
\end{tcolorbox}

\begin{tcolorbox}[knowledge,title=4.4 资产定价基本定理]
资产定价关系是从支付 $x$ 到价格的映射。记\textbf{定价算子}（估值算子）$V$：
\[
    S=V(x).
    \tag{4.1}
\]
严格地说，先在\textbf{可达支付空间} $\mathcal M=\{X\theta:\theta\in\mathrm R^N\}$ 上定义 $V(X\theta)=S^{\mathsf T}\theta$。若同一支付的不同复制组合成本不同，这个映射就没有唯一的值；一价定律正是保证它定义良好的条件。第 $n$ 只证券的价格即 $S_n=V(X_n)$。无套利把 $V$ 收成正的线性泛函。

\begin{tcolorbox}[knowledge]
\textbf{定理 4.3（一价定律）}
\[
    x=y\implies V(x)=V(y).
    \tag{4.2}
\]
\end{tcolorbox}

没有第 1 类套利就得到一价定律。推论：未来支付为 $0$ 的证券价格为 $0$，即 $V(0)=0$。证明：若 $x=y$ 却 $V(x)>V(y)$，卖出 $x$、买入 $y$，未来净支付为零而初始净收入为正，构成第 1 类套利。一价定律弱于无套利：没有第 1 类套利推不出没有第 2、3 类。

\begin{tcolorbox}[knowledge]
\textbf{定理 4.4}\quad 支付为正的证券价格为正：
\[
    x>0\implies V(x)>0.
    \tag{4.3}
\]

\textbf{定理 4.5}\quad $V$ 递增：
\[
    x_1\ge x_2\implies V(x_1)\ge V(x_2).
    \tag{4.4}
\]
\end{tcolorbox}

若 $x>0$ 而 $V(x)\le0$，买入它即为第 2 或第 3 类套利，故必须 $V(x)>0$。若 $x_1\ge x_2$ 却 $V(x_1)<V(x_2)$，买入 $x_1$、卖出 $x_2$，初始净收入为正且未来支付非负，也构成套利。当 $x_1-x_2>0$ 时还有严格结论 $V(x_1)>V(x_2)$。

\begin{tcolorbox}[knowledge]
\textbf{定理 4.6}\quad 无摩擦且无套利市场中，可达支付上的定价算子是递增的线性算子：对任意 $a,b\in\mathrm{R}$ 以及支付 $x$、$y$，
\[
    V(ax+by)=aV(x)+bV(y).
    \tag{4.5}
\]
从而 $V(0)=0$，并且存在向量 $\phi\in\mathrm{R}^{\Omega}$ 使 $V(x)=\phi^{\mathsf T}x$。
\end{tcolorbox}

\textbf{线性的证明：}分别复制 $x$、$y$，再持有 $a$ 倍与 $b$ 倍组合，得到支付 $ax+by$，成本为 $aV(x)+bV(y)$；一价定律要求它等于 $V(ax+by)$。允许负系数用到了卖空不受限的无摩擦假设。

这个向量表示首先只要求在 $\mathcal M$ 上成立；若市场不完全，不能据此为所有未交易支付直接指定唯一价格。形式与第 3 章的状态价格公式相同。这里只知道 $\phi$ 是一个向量；下一条才把它收成严格正、并且说明这正是无套利。

\begin{tcolorbox}[knowledge]
\textbf{定理 4.7（资产定价基本定理）}\quad 市场不存在套利机会，当且仅当存在 $\phi\gg 0$，使得每只可交易证券的价格都是
\[
    S=X^{\mathsf T}\phi.
    \tag{4.6}
\]
教材写成 $S=(\phi^{\mathsf T}X)^{\mathsf T}$，与上式相同。
\end{tcolorbox}

\textbf{充分性：}若 $S=X^{\mathsf T}\phi$、$\phi\gg0$，则任意组合满足 $S^{\mathsf T}\theta=\phi^{\mathsf T}X\theta$。未来支付为零时成本必为零；未来支付非负、非零时成本必严格为正。因此不可能同时有非正成本与套利支付，三类套利都被排除。必要性是 Stiemke 引理的特例：对 $m\times n$ 矩阵 $X$，存在 $\phi\gg 0$ 使 $S=X^{\mathsf T}\phi$，当且仅当集合 $\{\theta:[-S^{\mathsf T}\theta;\,X\theta]>0\}$ 为空。该集合正是套利机会的集合。证明放在教材附录。也可把扩展支付矩阵写成 $A=[-S^{\mathsf T};X]$：Stiemke 二择一定理说，“存在 $\theta$ 使 $A\theta>0$”与“存在 $y\gg0$ 使 $A^{\mathsf T}y=0$”恰有一个成立。无套利排除前者；把后者写成 $y=[y_0;y_1]$，则 $-y_0S+X^{\mathsf T}y_1=0$，取 $\phi=y_1/y_0\gg0$ 即得定价公式。

因此无套利意味着存在严格正的状态价格，可为\textbf{所有}可交易证券定价。这样的 $\phi$ 一般不唯一；定理 4.7 没有给出唯一性。给已有证券定价的线性解可以含负分量（练习 4.5），但那一支不能用来代表无套利：无套利要求的是存在一支 $\phi\gg 0$。

\begin{tcolorbox}[example]
\textbf{例 4.2}\quad 市场上先只有证券 $A$，价格为 $1$，三个状态下支付都是 $1$。没有套利。任一满足 $\phi_1+\phi_2+\phi_3=1$ 的 $\phi\gg 0$ 都能为它定价：
\[
    \Phi_A=\{\phi\gg 0:\phi_1+\phi_2+\phi_3=1\}.
\]
再加入证券 $B$，价格为 $1$，支付为 $[0;\,2;\,2]$。能为它定价的是
\[
    \Phi_B=\{\phi\gg 0:2\phi_2+2\phi_3=1\}.
\]
同时给 $A$、$B$ 定价的状态价格是交集
\[
    \Phi_A\cap\Phi_B=\{\phi\gg 0:\phi_1=\tfrac12,\;\phi_2+\phi_3=\tfrac12\},
\]
解不唯一，故 $A$、$B$ 的价格里没有套利。

再加入证券 $C$，价格为 $2$，支付为 $[2;\,0;\,0]$。条件变成 $\phi_1=1$，
\[
    \Phi_C=\{\phi\gg 0:\phi_1=1\}.
\]
三者的交集是空集。不存在能为所有证券定价的 $\phi\gg 0$，与例 4.1 中这三只证券存在套利一致。
\end{tcolorbox}

\begin{tcolorbox}[knowledge]
\textbf{定理 4.8}\quad 无套利的完全市场上，状态价格向量唯一。
\end{tcolorbox}

设 $\theta_\omega$ 复制状态 $\omega$ 的阿罗证券，$X\theta_\omega=\mathbf{1}_\omega$。则 $\phi_\omega=S^{\mathsf T}\theta_\omega$，由市场价格唯一给定；即使存在冗余证券而复制组合不唯一，一价定律仍保证所有复制组合成本相同。若有两支不同的定价向量 $\phi,\psi$，必有状态 $\omega$ 满足 $\phi_\omega\ne\psi_\omega$；但该阿罗证券已可复制，两支向量却给同一复制支付不同价格，与一价定律矛盾。方阵可逆时，直接有 $\phi=(X^{\mathsf T})^{-1}S$。完全意味着任意支付都可通过市场交易复制，因而其价格唯一；不完全市场的已有证券仍可有唯一市场报价，只是不能据此唯一确定所有未交易支付的价格。
\end{tcolorbox}

\begin{tcolorbox}[knowledge,title=4.5 风险中性定价和鞅]
设可交易证券中有一只无风险债券，1 期支付为 $\iota=[1,\ldots,1]^{\mathsf T}$，价格为 $S_1$。由基本定理，存在 $\phi\gg 0$ 使
\[
    S_1=\phi^{\mathsf T}\iota=\sum_{\omega\in\Omega}\phi_{\omega}.
\]
$1$ 单位债券投资的净收益率称为\textbf{无风险利率} $r_F$：
\[
    S_1(1+r_F)=1,\qquad r_F=\frac{1-S_1}{S_1}.
\]
因而
\[
    \frac{1}{1+r_F}=\sum_{\omega\in\Omega}\phi_{\omega}.
    \tag{4.7}
\]
债券价格可以写成对其确定支付 $1$ 的贴现。$1/(1+r_F)$ 是折现因子，$r_F$ 是折现率。贴现只是用状态价格定价的一种写法。利率也是货币的时间价值：今天 $1$ 单位资源换取未来确定资源时，市场给出的回报。

\begin{tcolorbox}[example,title=PPT 补充：从状态价格读出利率]
晴、雨两状态价格为 $\phi=[0.3,0.5]^{\mathsf T}$ 时，单位无风险债券价格为 $B=0.3+0.5=0.8$，所以 $r_F=1/0.8-1=25\%$。状态价格不是概率，其和不必为 $1$；规范化后 $Q=[3/8,5/8]$。总和 $0.8$ 决定时间贴现，各状态的相对权重决定不同风险支付如何估值，利率一个数并没有概括整个状态价格向量。
\end{tcolorbox}

对其余证券 $n=2,\ldots,N$，支付 $x_n=[x_{1n},\ldots,x_{\Omega n}]^{\mathsf T}$，
\[
    S_n=\phi^{\mathsf T}x_n=\sum_{\omega\in\Omega}\phi_{\omega}x_{\omega n}.
\]
定义\textbf{风险中性概率}（规范化的状态价格）
\[
    q_{\omega}=\frac{\phi_{\omega}}{\sum_{\omega'}\phi_{\omega'}}.
    \tag{4.8}
\]
$q_{\omega}>0$ 且 $\sum q_{\omega}=1$，故 $Q=\{q_{\omega}\}$ 是 $\Omega$ 上的概率测度。定价公式改写成
\[
    S_n=S_1\sum_{\omega}q_{\omega}x_{\omega n},
    \qquad
    S_n=\frac{\mathrm{E}^{Q}[\tilde{x}_n]}{1+r_F},\quad n=1,\ldots,N.
    \tag{4.9}
\]
(4.9) 就是 (4.6) 的另一种写法，称为\textbf{风险中性定价}（risk-neutral pricing），$Q$ 称为\textbf{风险中性测度}。它适用于所有证券，并不区分风险大小：在 $Q$ 下，每只证券的价格都只是期望支付按无风险利率折现，好像只取得无风险收益。

\textbf{为什么不需要假设投资者真的风险中性？}风险偏好、禀赋等会影响基础证券的均衡价格；给定这些价格后，可复制支付的价格由无套利约束决定，不必重新询问投资者偏好。$Q$ 是把已有价格中的风险调整吸收到权重后的定价工具。若市场不完全，已有证券价格通常只约束一组可行 $Q$，无法由此给所有衍生支付一个唯一价格。

\begin{tcolorbox}[example,title=PPT 补充：套餐支付的价格由复制成本决定]
股票 $A$ 价格为 $10$，两状态支付为 $[20,0]^{\mathsf T}$；债券 $B$ 价格为 $5$，支付为 $[5,5]^{\mathsf T}$。衍生证券 $C$ 支付为 $[25,5]^{\mathsf T}$，等于 $A+B$，所以价格必须为 $15$。若报价为 $20$，卖出 $C$ 并买入 $A+B$，初始收入 $5$，未来净支付为零。

由 $20\phi_1=10$、$5(\phi_1+\phi_2)=5$，得 $\phi_1=\phi_2=1/2$，无风险利率为零，故 $C=\mathrm E^Q[\tilde C]=(25+5)/2=15$。复制法与风险中性法给出同一价格；这里 $Q$ 两状态相等，不代表真实概率也相等。
\end{tcolorbox}

$Q$ 不必等于真实概率 $P$。它由 (4.8) 从状态价格定义，风险已经折进 $q_{\omega}$。因此 (4.9) 是定价关系，不是用 $P$ 作的概率预期。

\begin{tcolorbox}[example]
\textbf{例 4.3}\quad 1 期有两个等可能状态 $a$、$b$。
\begin{center}
\begin{tabular}{@{}c@{\quad}c@{}}
\textbf{证券 1} & \textbf{证券 2} \\[0.35em]
\WJEventTree[left={$1$},stem]{1,1} &
\WJEventTree[left={$\tfrac12$},stem]{2,0}
\end{tabular}
\end{center}
证券 1 是无风险债券，$r_F=0$。在真实概率下两只证券的期望支付都是 $1$，但证券 2 的价格只有证券 1 的一半：期望值相同，确定支付与风险支付的价格不必相同。

由价格得 $\phi_a+\phi_b=1$、$2\phi_a=1/2$，故 $\phi_a=1/4$、$\phi_b=3/4$，$Q=\{1/4,\,3/4\}$。风险中性定价给出
\[
    \frac{\mathrm{E}^{Q}[\tilde{x}_1]}{1+0}
    =\tfrac14\cdot 1+\tfrac34\cdot 1=1,
    \qquad
    \frac{\mathrm{E}^{Q}[\tilde{x}_2]}{1+0}
    =\tfrac14\cdot 2+\tfrac34\cdot 0=\tfrac12,
\]
正是两只证券的价格。求期望时用的是 $Q$ 不是 $P$。
\end{tcolorbox}

因为 $S_1=1/(1+r_F)$ 且债券的 1 期支付 $\tilde{x}_1=1$，(4.9) 也可写成
\[
    \frac{S_n}{S_1}=\mathrm{E}^{Q}\Bigl[\frac{\tilde{x}_n}{\tilde{x}_1}\Bigr].
\]
$S_n$、$S_1$ 以现时消费品计价；$S_n/S_1$ 则是以无风险债券计价的价格。用债券做计价单位时，0 期价格是 $S_n/S_1$，1 期价格是 $\tilde{x}_n$（债券在 1 期的价格恒为 $1$）。这改变的是计价单位，使时间贴现体现在分母中；它并没有消除风险，风险调整仍包含在 $Q$ 中。记 $\hat{S}_{n,t}=S_{n,t}/S_{1,t}$（$t=0,1$），则
\[
    \hat{S}_{n,0}=\mathrm{E}^{Q}[\hat{S}_{n,1}].
    \tag{4.10}
\]

\begin{tcolorbox}[knowledge]
\textbf{定义 4.3}\quad 设 $\mathcal F_t$ 表示截至 $t$ 期已知的信息。随机过程 $z_t$ 若适应该信息、满足 $\mathrm E^Q|z_t|<\infty$，且
\[
 z_t=\mathrm E^Q[z_{t+1}\mid\mathcal F_t],
\]
则称为测度 $Q$ 下的\textbf{鞅}。当前信息下，下一期的预期变动为零；这不表示每次实现值都不变，也不表示真实概率下的预期收益为零。
\end{tcolorbox}

(4.10) 说的是：以债券计价后，证券价格在风险中性测度 $Q$ 下是鞅。因此 $Q$ 也称为\textbf{等价鞅测度}。这里的“等价”是指 $Q$ 与真实测度 $P$ 有相同的零概率集：$p_{\omega}=0$ 当且仅当 $q_{\omega}=0$。$\phi\gg 0$ 给出 $q\gg 0$；若真实概率也严格为正，二者等价。

\textbf{等价性的套利含义须落实到可交易支付。}若事件 $E$ 在真实世界可能发生，即 $P(E)>0$，而候选 $Q'$ 令 $Q'(E)=0$，那么按它给事件保险 $\mathbf1_E$ 定价会得到零价格。若该支付确实可交易（或可复制），免费买入便是第 2 类套利。\textbf{单纯“低估某项风险”或买入后可能获利，并不足以构成套利}；看跌期权也不一定只在破产时支付。若 $\mathbf1_E$ 在不完全市场中不可交易、不可复制，仅发现某候选定价测度漏掉 $E$，不能直接声称已有证券中存在套利，仍须寻找可行交易组合。

本书只有两个时点，式 (4.10) 只取 $t=0$。规范化价格服从鞅这一关系可以直接搬到动态框架；经济含义在这个静态模型里已经清楚。
\end{tcolorbox}

\begin{tcolorbox}[knowledge,title=本章小结]
\begin{enumerate}
    \item 无摩擦时，线性无关的证券组合给出市场结构的等价描述。市场完全当且仅当独立支付的证券数等于状态数。
    \item 套利是 $S^{\mathsf T}\theta\le 0$、$X\theta\ge 0$ 且至少一处严格。它只依赖公开的支付和价格，不依赖私有信息或真实概率；无摩擦时人人都可以做。
    \item 均衡中不存在套利。前提是市场无摩擦，且至少有人严格偏好非负、非零的消费增加。无套利是均衡的必要条件。
    \item 资产定价基本定理：无套利当且仅当存在 $\phi\gg 0$ 使 $S=X^{\mathsf T}\phi$。完全市场中这个 $\phi$ 唯一。
    \item 风险中性定价 $S_n=\mathrm{E}^{Q}[\tilde{x}_n]/(1+r_F)$ 是同一公式的改写。$Q$ 由状态价格规范化得到，不必等于真实概率。
    \item 用无风险债券计价后，价格在 $Q$ 下是鞅，$Q$ 是等价鞅测度。
\end{enumerate}
\end{tcolorbox}
\fi

% !TeX root = ../../金融经济学_王江_讲义及习题.tex
% 04_套利和资产定价 课后习题（由章节文件拆分）

\begin{problem}
\problemno{4.1} 经济在 1 期有 4 个可能状态。在市场中有 5 只可交易证券，它们的支付矩阵 $X$ 如下：
\[
    X=\begin{bmatrix}
        1 & 1 & 0 & 0 & 0 \\
        1 & 2 & 1 & 0 & 0 \\
        1 & 3 & 2 & 1 & 0 \\
        1 & 4 & 3 & 2 & 1
    \end{bmatrix}.
\]
\begin{enumerate}[label=(\alph*)]
    \item 证明市场是完全的。
    \item 假设五只证券的 0 期价格分别为 $1$，$2.5$，$1.5$，$0.75$，$0.25$。证明状态价格向量存在且唯一，并给出它的值。
    \item 证明市场上存在冗余证券。
    \item 选择一组足以保证市场完全的复合证券。把每一只 Arrow-Debreu 证券都表示成这些证券的组合。称这些组合为状态或有组合。
    \item 用状态或有组合把任一由五只交易证券组成的组合表示出来。
    \item 证明状态或有组合可生成状态空间。
\end{enumerate}
\end{problem}
\begin{solution}
记 5 只交易证券的支付向量分别为
\[
    X_{1}=\begin{bmatrix}1 \\ 1 \\ 1 \\ 1\end{bmatrix},\quad X_{2}=\begin{bmatrix}1 \\ 2 \\ 3 \\ 4\end{bmatrix},\quad X_{3}=\begin{bmatrix}0 \\ 1 \\ 2 \\ 3\end{bmatrix},\quad X_{4}=\begin{bmatrix}0 \\ 0 \\ 1 \\ 2\end{bmatrix},\quad X_{5}=\begin{bmatrix}0 \\ 0 \\ 0 \\ 1\end{bmatrix}
\]
\begin{enumerate}[label=(\alph*)]
    \item 因为 $X_{1}$、$X_{2}$、$X_{3}$、$X_{4}$ 的支付是互相独立的四只证券，而状态数等于 4，因而证券市场是完全的；
    \item 假设状态价格向量为 $\phi=\begin{bmatrix}\phi_{1} \\ \phi_{2} \\ \phi_{3} \\ \phi_{4}\end{bmatrix}$，那么有：
    \[
        \begin{aligned}
            \phi_{1}+\phi_{2}+\phi_{3}+\phi_{4}&=1\\
            \phi_{1}+2\phi_{2}+3\phi_{3}+4\phi_{4}&=2.5\\
            \phi_{2}+2\phi_{3}+3\phi_{4}&=1.5\\
            \phi_{3}+2\phi_{4}&=0.75\\
            \phi_{4}&=0.25
        \end{aligned}
    \]
    求解上述方程组可以得到唯一的解为 $\phi=\begin{bmatrix}\tfrac14 \\ \tfrac14 \\ \tfrac14 \\ \tfrac14\end{bmatrix}$；
    \item 买入 1 单位的证券 2 卖出 1 单位的证券 3，得到的支付是 $\begin{bmatrix}1 \\ 1 \\ 1 \\ 1\end{bmatrix}$，恰好就是证券 1 的支付，因而市场上存在冗余证券；
    \item 证券 2、3、4、5 可以保证市场完全，而且我们有
    \[
        \begin{aligned}
            X_{2}-2X_{3}+X_{4}&=\begin{bmatrix}1 \\ 0 \\ 0 \\ 0\end{bmatrix}\\
            X_{3}-2X_{4}+X_{5}&=\begin{bmatrix}0 \\ 1 \\ 0 \\ 0\end{bmatrix}\\
            X_{4}-2X_{5}&=\begin{bmatrix}0 \\ 0 \\ 1 \\ 0\end{bmatrix}\\
            X_{5}&=\begin{bmatrix}0 \\ 0 \\ 0 \\ 1\end{bmatrix}
        \end{aligned}
    \]
    \item 假设组合中 5 只交易证券的头寸分别为 $\theta_{1}$、$\theta_{2}$、$\theta_{3}$、$\theta_{4}$、$\theta_{5}$，那么状态或有组合的头寸分别为 $\theta_{1}+\theta_{2}$，$\theta_{1}+2\theta_{2}+\theta_{3}$，$\theta_{1}+3\theta_{2}+2\theta_{3}+\theta_{4}$，$\theta_{1}+4\theta_{2}+3\theta_{3}+2\theta_{4}+\theta_{5}$；
    \item 因为状态或有组合的支付是 $\mathrm{R}^{4}$ 空间的一组基，因而状态或有组合可生成支付空间。
\end{enumerate}
\end{solution}

\begin{problem}
\problemno{4.2} 银行 1 的存贷利率为 $r_{1}$ 而银行 2 的为 $r_{2}$，且 $r_{1}>r_{2}$。证明存在套利机会。
\end{problem}
\begin{solution}
从银行 2 借钱并全部存入银行 1，参与者就可以获得第 2 类套利。
\end{solution}

\begin{problem}
\problemno{4.3} 银行 $i$ 的贷款利率为 $r_{ia}$（即顾客借钱的利率）而存款利率为 $r_{ib}$，$i=1,2$。
\begin{enumerate}[label=(\alph*)]
    \item 证明无套利要求 $r_{ia}\ge r_{ib}$。
    \item 两个银行的存贷利率应该满足什么样的条件才能使市场上不存在套利机会。
\end{enumerate}
\end{problem}
\begin{solution}
\begin{enumerate}[label=(\alph*)]
    \item 设 $r_{ia}<r_{ib}$。从银行 $i$ 借入 $1$ 并存入该银行：0 期收到贷款 $1$、支付存款 $1$，净现金流为 $0$；1 期收回 $1+r_{ib}$、偿还 $1+r_{ia}$，净现金流为
    \[
        (1+r_{ib})-(1+r_{ia})=r_{ib}-r_{ia}>0.
    \]
    这是第 2 类套利。因此无套利要求 $r_{ia}\ge r_{ib}$。
    \item 无套利当且仅当最高存款利率不超过最低贷款利率：
    \[
        \max(r_{1b},r_{2b})\le\min(r_{1a},r_{2a}),
    \]
    即 $r_{1a}\ge r_{1b}$、$r_{2a}\ge r_{2b}$、$r_{1a}\ge r_{2b}$、$r_{2a}\ge r_{1b}$。

    若 $r_{1a}<r_{2b}$，从银行 1 借入 $1$ 并存入银行 2，0 期净现金流为 $0$，1 期净现金流为 $r_{2b}-r_{1a}>0$。若 $r_{2a}<r_{1b}$，从银行 2 借入并存入银行 1，1 期净现金流为 $r_{1b}-r_{2a}>0$。本行的情形已由 (a) 给出。

    反过来，设上述四个不等式成立，记 $r_{b}^{*}=\max(r_{1b},r_{2b})$、$r_{a}^{*}=\min(r_{1a},r_{2a})$，则 $r_{b}^{*}\le r_{a}^{*}$。以 $\theta_{i}\ge 0$ 表示在银行 $i$ 的借款、$\delta_{i}\ge 0$ 表示存款。0 期净流入为 $\sum\theta_{i}-\sum\delta_{i}$，1 期净流入为
    \[
        \sum_{i}\delta_{i}(1+r_{ib})-\sum_{i}\theta_{i}(1+r_{ia})
        \le(1+r_{b}^{*})\sum\delta_{i}-(1+r_{a}^{*})\sum\theta_{i}.
    \]
    若 0 期净流入为 $0$，则 $\sum\theta_{i}=\sum\delta_{i}$，1 期净流入不超过 $(r_{b}^{*}-r_{a}^{*})\sum\delta_{i}\le 0$。若 0 期净流入为正，则 $\sum\theta_{i}>\sum\delta_{i}$，1 期净流入严格为负。因此不存在套利。
\end{enumerate}
\end{solution}

\begin{problem}
\problemno{4.4} 考虑练习 4.1 中定义的市场。假设 5 只交易证券的价格分别为 $1$，$2$，$3/2$，$1$，$1/2$。
\begin{enumerate}[label=(\alph*)]
    \item 在这样的价格和支付下存在套利机会吗？
    \item 你能构建第 1、2 以及 3 类套利机会吗？
\end{enumerate}
\end{problem}
\begin{solution}
\begin{enumerate}[label=(\alph*)]
    \item 买入 1 单位的证券 4，卖出 2 单位的证券 5，0 期支付为 0，而 1 期支付为 $\begin{bmatrix}0 \\ 0 \\ 1 \\ 0\end{bmatrix}>0$，因而存在套利机会。
    \item 上面构造了第 2 类套利机会。卖出 1 单位的证券 1 和 1 单位的证券 3，同时买入 1 单位的证券 2，那么 0 期支付为 $\tfrac12$，而 1 期支付为 0，这是第 1 类套利机会。同时进行上述两种交易就可以得到第 3 类套利机会。
\end{enumerate}
\end{solution}

\begin{problem}
\problemno{4.5} 在 1 期有三个可能状态：$\{\omega_{1},\omega_{2},\omega_{3}\}$。在经济中有三只交易证券，$a$，$b$，$c$，它们的支付向量如下：
\[
    X_{a}=\begin{bmatrix}1\\1\\0\end{bmatrix},\quad
    X_{b}=\begin{bmatrix}0\\1\\1\end{bmatrix},\quad
    X_{c}=\begin{bmatrix}1/2\\1\\1/2\end{bmatrix}.
\]
这些证券的 0 期均衡价格为：$S_{a}=S_{b}=S_{c}=1$。令 $\phi=\begin{bmatrix}\phi_{1} \\ \phi_{2} \\ \phi_{3}\end{bmatrix}$ 为状态价格向量。
\begin{enumerate}[label=(\alph*)]
    \item 证明 $\begin{bmatrix}2 \\ -1 \\ 2\end{bmatrix}$ 是一个可能的状态价格向量，因为它给出了三支现有证券的价格。
    \item “资产定价基本定理”指出状态价格必须是非负的。在一个两期经济里，对于证券市场这意味着什么？并给出其解释。
    \item 上面的例子是这个原理的反例吗？为什么？
    \item 证明 $\begin{bmatrix}1/2 \\ 1/2 \\ 1/2\end{bmatrix}$ 是另外一个可能的状态价格向量。在这种情况下，所有状态价格都是正的。
    \item 解释为什么存在多个可能的状态价格向量，它们对已有证券都给出相同的价格。
\end{enumerate}
\end{problem}
\begin{solution}
\begin{enumerate}[label=(\alph*)]
    \item 因为
    \[
        2\times 1+(-1)\times 1=1,\quad (-1)\times 1+2\times 1=1,\quad 2\times\tfrac12+(-1)\times 1+2\times\tfrac12=1
    \]
    因而 $\begin{bmatrix}2 \\ -1 \\ 2\end{bmatrix}$ 是一个可能的状态价格向量；
    \item 资产定价基本定理：没有套利，当且仅当存在严格正的状态价格 $\phi\gg 0$，使每只可交易证券的价格都是 $S=X^{\mathsf T}\phi$。因此“状态价格必须非负”在这里的含义是均衡价格里不能有套利：若每一支给现有证券定价的 $\phi$ 都有负分量，市场就存在套利。
    \item 不是。因而它只是可能的状态价格之一，并不一定是真实的状态价格；定理要求的是存在一支 $\phi\gg 0$，而不是每一个解都严格为正。下面 (d) 给出的正向量说明无套利仍然成立；
    \item 因为
    \[
        \tfrac12\times 1+\tfrac12\times 1=1,\quad \tfrac12\times 1+\tfrac12\times 1=1,\quad \tfrac12\times\tfrac12+\tfrac12\times 1+\tfrac12\times\tfrac12=1
    \]
    因而 $\begin{bmatrix}\tfrac12 \\ \tfrac12 \\ \tfrac12\end{bmatrix}$ 是一个可能的状态价格向量；
    \item 因为市场是不完全的，因而可能存在多个状态价格，它们都给出已有证券的价格。
\end{enumerate}
\end{solution}

\begin{problem}
\problemno{4.6} 远期合约（Forward）是在未来（即 1 期）以当期（即 0 期）确定的价格交易某一证券或商品的合约，而在当期没有现金流产生。在当期约定的未来交易价格称作远期价格。比方说，考虑一份在 1 期以 $F_{0}$ 的价格购买 1 盎司黄金的远期合约。合约的买方同意在 1 期（合约的到期日）从卖方买入 1 盎司黄金。在当期没有现金交换。令 $S_{t}$ 为 $t=0,1$ 期黄金的价格，而 $r_{F}$ 为利率。考虑在 1 期得到 1 盎司黄金的两种方法：
\begin{itemize}
    \item 现在购买一份远期合约。在 1 期，支付远期价格 $F_{0}$ 并获得黄金。
    \item 现在支付 $S_{0}$ 买入 1 盎司黄金，并且持有它直到 1 期。假设从现在起持有黄金直到 1 期没有另外的成本或收益。
\end{itemize}
在这两种情况下，我们在 1 期同样得到 1 盎司黄金，而那时它的价值为 $S_{1}$。
\begin{enumerate}[label=(\alph*)]
    \item 由于这两种策略给予我们相同的支付 $S_{1}$，无套利原理意味着它们的成本必须一样。用当期的价格 $S_{1}$ 和 $r_{F}$ 给出远期价格。
    \item 假设持有 1 盎司黄金不仅得到它的 1 期出售收益，在 1 期还得到额外收益 $y$。$y$ 将如何影响远期价格？
\end{enumerate}
\end{problem}
\begin{solution}
\begin{enumerate}[label=(\alph*)]
    \item 第一种方法的成本是 $\dfrac{F_{0}}{1+r_{F}}$，而第二种方法的成本是 $S_{0}$，根据无套利原理，$\dfrac{F_{0}}{1+r_{F}}=S_{0}$，因此，$F_{0}=(1+r_{F})S_{0}$；
    \item 由于第二种方法还得到额外收益，由无套利原理原理有 $F_{0}<(1+r_{F})S_{0}$。（如果收益是确定性的，那么应该有 $F_{0}+y=(1+r_{F})S_{0}$）
\end{enumerate}
\end{solution}

\begin{problem}
\problemno{4.7} 英镑现时对美元的价格（即英镑/美元汇率）为 $S_{0}$（即一英镑的美元价格为本 $S_{0}$），美元利率为 $r_{A}$ 而英镑利率为 $r_{B}$。考虑一份在 1 期购买 1 英镑的远期合约。远期价格（即远期汇率）应该是多少？提示：考虑在 1 期得到 1 英镑的两种策略。
\begin{itemize}
    \item 购买一份远期合约。在 1 期，支付远期价格 $F_{0}$ 购入 1 英镑。
    \item 现在以当前汇率 $S_{0}$ 买入 $x$ 英镑并以 $r_{B}$ 的利率进行投资，在 1 期得到 1 英镑。
\end{itemize}
得到的远期价格、现时汇率和利率之间的关系就叫作抛补利率平价理论（Covered Interest Rate Parity）。
\end{problem}
\begin{solution}
第一种策略的 0 期支付是 $\dfrac{F_{0}}{1+r_{A}}$ 美元，而第二种策略的 0 期支付是 $\dfrac{S_{0}}{1+r_{B}}$ 美元，根据无套利原理，$\dfrac{F_{0}}{1+r_{A}}=\dfrac{S_{0}}{1+r_{B}}$，即 $F_{0}=\dfrac{1+r_{A}}{1+r_{B}}S_{0}$。
\end{solution}

\begin{problem}
\problemno{4.8} 考虑练习 4.1 定义的市场。假设证券 1，2，4，5 的价格分别是 $0.8$，$2$，$0.6$，$0.2$。
\begin{enumerate}[label=(\alph*)]
    \item 计算状态价格和利率。
    \item 构建等价鞅测度。
    \item 证明 5 只证券的价格都具有鞅性质。
\end{enumerate}
\end{problem}
\begin{solution}
记练习 4.1 中五只证券的支付列向量为 $X_{1},\ldots,X_{5}$。状态价格 $\phi=\begin{bmatrix}\phi_{1} \\ \phi_{2} \\ \phi_{3} \\ \phi_{4}\end{bmatrix}$ 满足 $S_{n}=\phi^{\mathsf T}X_{n}$，即
\[
    S_{1}=\phi^{\mathsf T}X_{1},\quad S_{2}=\phi^{\mathsf T}X_{2},\quad S_{4}=\phi^{\mathsf T}X_{4},\quad S_{5}=\phi^{\mathsf T}X_{5}.
\]
\begin{enumerate}[label=(\alph*)]
    \item 代入 $S_{1}=0.8$，$S_{2}=2$，$S_{4}=0.6$，$S_{5}=0.2$ 得
    \[
        \begin{aligned}
            \phi_{1}+\phi_{2}+\phi_{3}+\phi_{4}&=0.8,\\
            \phi_{1}+2\phi_{2}+3\phi_{3}+4\phi_{4}&=2,\\
            \phi_{3}+2\phi_{4}&=0.6,\\
            \phi_{4}&=0.2.
        \end{aligned}
    \]
    由后两式 $\phi_{3}=0.2$；前两式相减得 $\phi_{2}=0.2$，$\phi_{1}=0.2$。故
    \[
        \phi=\begin{bmatrix}0.2 \\ 0.2 \\ 0.2 \\ 0.2\end{bmatrix},\qquad
        S_{3}=\phi^{\mathsf T}X_{3}=1.2.
    \]
    无风险利率 $r_{F}=\bigl(\sum_{i}\phi_{i}\bigr)^{-1}-1=0.25$。
    \item 等价鞅测度 $Q_{i}=\phi_{i}/\sum_{j}\phi_{j}$，即 $Q=\begin{bmatrix}0.25 \\ 0.25 \\ 0.25 \\ 0.25\end{bmatrix}$。
    \item 两期模型里的鞅性质是风险中性定价：用等价鞅测度 $Q$ 对 1 期支付取期望，再按无风险利率折现，应当回到 0 期价格，即 $S_{n}=\mathrm{E}^{Q}[X_{n}]/(1+r_{F})$。由 $Q_{i}=\phi_{i}/\sum_{j}\phi_{j}$ 与 $1+r_{F}=1/\sum_{j}\phi_{j}$，
    \[
        \frac{\mathrm{E}^{Q}[X_{n}]}{1+r_{F}}=\phi^{\mathsf T}X_{n}.
    \]
    证券 $1$、$2$、$4$、$5$ 的价格正是用来解出 $\phi$ 的，故 $\phi^{\mathsf T}X_{n}=S_{n}$。证券 $3$ 的价格取为 $S_{3}=\phi^{\mathsf T}X_{3}$。在 $Q_{i}=1/4$、$1+r_{F}=5/4$ 下，
    \[
        \begin{aligned}
            \frac{\mathrm{E}^{Q}[X_{1}]}{1+r_{F}}&=\frac{(1+1+1+1)/4}{5/4}=0.8=S_{1},\\
            \frac{\mathrm{E}^{Q}[X_{2}]}{1+r_{F}}&=\frac{(1+2+3+4)/4}{5/4}=2=S_{2},\\
            \frac{\mathrm{E}^{Q}[X_{3}]}{1+r_{F}}&=\frac{(0+1+2+3)/4}{5/4}=1.2=S_{3},\\
            \frac{\mathrm{E}^{Q}[X_{4}]}{1+r_{F}}&=\frac{(0+0+1+2)/4}{5/4}=0.6=S_{4},\\
            \frac{\mathrm{E}^{Q}[X_{5}]}{1+r_{F}}&=\frac{(0+0+0+1)/4}{5/4}=0.2=S_{5}.
        \end{aligned}
    \]
\end{enumerate}
\end{solution}


